
%\needspace{4\baselineskip}
\color{uniblue}{\subsection{ Автоматический ввод резерва (АВР)}}
\color{black}

\begin{enumerate}

%%%%%%%%%%%%%%%%%%%%%%%%%%%%%%%%%%%%%%%%%%%%%%%%%%%%%%%%%%%%%%%%%%%%%%  ОБЩИЕ СВЕДЕНИЯ %%%%%%%%%%%%%%%%%%%%%%%%%%%%%%%%%%%%%%%%%%%%%%%%%%

\item Описание функции АВР

\begin{enumerate}

\item
В устройстве автоматический ввод резерва представлен функцией «АВР». Функция  предназначена  для  восстановления  питания  потребителей  путем  автоматического присоединения резервного источника.

\item
В устройстве реализовано два режима работы АВР определяемые уставкой SGF2 <<Реж\_АВР>>:
\begin{itemize} 
    \item внутренний; 
    \item внешний. 
\end{itemize}

\item
Описание алгоритмов работы АВР СВ:
\begin{enumerate}[label=\arabic*), noitemsep]
    \item Режим работы <<Внутренний>>:
    \begin{itemize}
        \item  В  этом  режиме  при  наличии  отключенного  положения  СВ  <<В: Отключен>>,  отсутствии  блокировок  и  наличие  контроля  напряжения  на  1  и  2  секции  шин  (контроль определяется уставкой SGF2) по истечении времени таймера Т2 на входе <<S>> RS–триггера формируется <<1>> на время таймера Т4 (5 мс). В результате чего при отсутствии <<1>> на входе <<R>> формируется сигнал готовности АВР <<АВР: Готовность>> для включения вводного выключателя 1 или 2 секции шин. 

        \item С  помощью  уставки  SGF3 <<Реж\_КОН>> контролируется  отсутствие  напряжения  на  секции  шин.  В  момент пропадания  напряжения  на  1  секции  шин  формируется  выходной  сигнал  разрешения  на  отключение вводного выключателя 1 секции шин <<АВР: Разреш откл ВВ 1сш>>. 

        \item В  данном  условии  присутствует  блокировка  от  вводного  выключателя  2  секции  шин,  что предотвращает  одновременное  питание  нагрузки  от  двух  источников,  что  может  привести  к  аварийной ситуации. 

        \item Выходной сигнал <<АВР: Разреш откл ВВ 1сш>> направляется в терминал вводного выключателя 1 секции шин. После отключения вводного выключателя 1 секции шин. на дискретный вход терминала СВ 
        поступает  сигнал  <<ВВ  1сш  отключен>>.  При  наличии  данных  условий  на  время  таймера  Т3, продолжительность которого составляет 10 секунд, формируется выходной сигнал пуска АВР <<АВР: Пуск>>.

        \item При  наличии  пусковых  условий  запускается  таймер  Т1  и  по  истечении  его  времени  проверяется наличие блокировки включения из логики КСВ. Если блокировки нет, то формируется команда на включение секционного выключателя <<АВР: Включение>> и сброс готовности схемы. 

        \item При  отсутствии готовности схемы АВР,  таймер  Т3,  который  отвечает  за  время  ввода  АВР  перед формированием  команды  на  включение  резервного  источника  питания,  сбрасывается.  Таким  образом, процесс АВР не  продолжается,  пока  не  будет  получен  сигнал  готовности  схемы.  Это  помогает предотвратить ошибочное включение резервного питания в случае, если система не готова к этому.

        \item Отключение вводного выключателя второй секции шин можно выполнить по аналогии с алгоритмом для первой секции. Отключение вводного выключателя второй секции шин может быть выполнено в соответствии с алгоритмом, аналогичным тому, который используется для первой секции.
    \end{itemize}

    \item Режим работы <<Внешний>>:
    \begin{itemize}
        \item  В этом режиме контроль напряжения осуществляется на вводных выключателях 1 и 2 секции шин. Это значит, что основное управление и мониторинг напряжения осуществляется от терминалов защит вводных выключателей.
        \item Если на терминал СВ поступает сигнал <<Вкл СВ от ВВ 1сш>> или <<Вкл СВ от ВВ 2сш>>, это указывает на то, что вводные выключатели секций шин готовы к включению секционного выключателя. При наличии готовности АВР и отсутствии сигнала блокировки включения из логики КСВ, команда на включение СВ и сброс готовности схемы АВР формируются без выдержки времени. 
    \end{itemize}
\end{enumerate}

\item
Поясняющие схемы организации цепей АВР и ВНР показаны на рисунках \ref{appavr:img1} и \ref{appavr:img2} приложения \ref{app:avr}. 

\item
Ввод АВР в работу на этапе параметрирования устройства осуществляется программным переключателем SGF1 <<Ввод\_функции>>. Информация о введенной функции отображается выходным сигналом <<АВР: Ввод>>.
\item
Функция может быть выведена из работы оперативно путем активации сигнала <<ОВ АВР>>. Выведенное состояние функции характеризуется активным состоянием сигнала <<АВР: Оперативный вывод>>.

\item
Алгоритм работы функции выполнен в соответствии с рисунком \ref{avr:img1}. В таблице \ref{avr:tbl1} приведены параметры, необходимые для настройки функции.

\vspace{3mm}
\begin{figure}[H]
\centering
\includegraphics[width=0.9\textwidth,height=0.9\textheight,keepaspectratio]{\fbpath/410. АВР СВ/_latex/img/avr.pdf}
\captionof{figure}{Функциональная схема логики АВР}\label{avr:img1}
\end{figure}

\end{enumerate}

%%%%%%%%%%%%%%%%%%%%%%%%%%%%%%%%%%%%%%%%%%%%%%%%%%%%%%%%%%%%%%%%%%%%%% ЗАПРЕТ АВР %%%%%%%%%%%%%%%%%%%%%%%%%%%%%%%%%%%%%%%%%%%%%%%%%%
\item Формирование запрета АВР (ЗАВР)
\begin{enumerate}

\item
В устройстве запрет автоматического ввода резерва представлен функцией «ЗАВР». Запрет АВР позволяет избежать непреднамеренного включения оборудования, что может привести к повреждению или нарушению безопасности.

\item
Запрет АВР активируется, когда введена функция АВР и поступил один из сигналов как из логики, так и на дискретные входы устройства.

\item
Запрет АВР от устройства ТН функции сигнализации замыкания на землю (СЗЗ) первой или второй секции шин контролируется уставкой SGF1. 

\item
Алгоритм работы функции выполнен в соответствии с рисунком \ref{zavr:img1}. В таблице \ref{avr:tbl1} приведены параметры, необходимые для настройки функции.

\vspace{3mm}
\begin{figure}[H]
\centering
\includegraphics[width=0.85\textwidth,height=0.85\textheight,keepaspectratio]{\fbpath/410. АВР СВ/_latex/img/zavr.pdf}
\captionof{figure}{Функциональная схема логики ЗАВР}\label{zavr:img1}
\end{figure}

\end{enumerate}

\end{enumerate}
